\chapter[Sermon 15. He Blames Certain Things in the Congregation of Sardis: Notwithstanding He Shows Straightway a Remedy, Whereby They May Be Healed, and Be Safe.]{He Blames Certain Things in the Congregation of Sardis: Notwithstanding He Shows Straightway a Remedy, Whereby They May Be Healed, and Be Safe.}
\label{sermon:015}
\markright{Sermon 15. He Blames Certain Things in the Congregation of ...}

And write to the messenger of the congregation of Sardis, this says he who has the spirits of God and the seven stars. I know your works: you have a reputation that you live, but you are dead. Be watchful and strengthen what remains, which is ready to die. For I have not found your works perfect before God. Remember therefore how you have received and heard, and hold fast and repent. If you do not watch, I will come on you as a thief, and you will not know the hour when I will come upon you.

In one congregation of Sardis there were two sorts of people, professing on either side the name of Christ. But some in deed answered but little to the holy profession, living more licentiously than became them. And the others in holiness of life set forth the doctrine of our Savior that they professed. The Lord Jesus accuses the first sort in this Epistle by John, and shows also a medicine for the disease. And he exhorts the latter to perseverance, commending their integrity. Therefore this Epistle is divided into two parts, very fit and profitable for our time.

The first part of the Epistle contains those things which we have now recited. He does not proceed in this matter in any other order than we have seen him proceed hitherto. For first, he shows to whom it is dedicated and sent: namely, to the pastor of the congregation of Sardis, and therefore also to the whole church. Sardis is said to have been the head city of Lydia or of Maonia, the metropolitan city of Croesus, the most rich king of Lydia, whom Herodotus writes that King Cyrus overcame, a town most famous, and steeped and stained with pride, that it was a den of vice and addicted to voluptuousness. For Strabo in the thirteenth book of Geography testifies that all the maidens thereof were harlots, who mentions more of the same city. Certainly, it seems to have kept its old ways even at such a time as it had received the name of the Lord, and therefore to have been more given to fornication and all manner of filthy lust. The Lord seems to have blamed this in them, as Saint Paul likewise persecuted the same vice in the Corinthians. The world can hardly believe that simple fornication is a sin, whereupon in that great council of the Apostles, both they and the elders and the whole assembly with one mind decreed that the Gentiles should abstain from fornication. The devil at this day goes about many times to defile the church again with fornication, to set up brothels, and that by authority and openly whoredom might be practiced. For so being cast out, he takes seven worse spirits, enterprising to possess that place again out of which he was exiled by the preaching of the Gospel. We must therefore resist him, lest the Lord Jesus himself do accuse us, as he does here accuse them of Sardis most grievously. Then the Lord Jesus is declared to be author of the Epistle, not without praise. For he is said to have the seven spirits of God, that is to have the seven formed spirit, whom he also pours out upon the faithful, or else he is one only spirit, and not seven: but seven, that is to wit, his graces are many and diverse, as I declared in the first chapter. For he also has in his right hand seven stars, to wit, the whole multitude of all preachers and ministers, keeping and instructing them. And this beginning agrees not amiss with this argument, which he treats in this Epistle. For of the spirit of Christ is life; of the want of the spirit, death. Christ preserves the ministers, however angry some are in the church, for accusing their wickedness. Privily therefore he warns them to crave the spirit, to nourish the spiritual life, and to trust in Christ, who will defend the ministers and advance them.

\index{John, Saint}As he testifies in all his other letters, he repeats it here also: “I know your works.” Of this I have spoken before. The Lord is ignorant of nothing that is done in the church, which is also the searcher of hearts. And especially he condemns this in this church, that she thought herself alive where she was dead. He speaks not of physical, but of spiritual life and death. For Christ lives by his spirit in his saints and faithful, and shows lively works through them, as the Lord teaches in John 6 and many other places in the Gospel of Saint John. The Apostle said also that he did not now live, but that Christ lived in him. The same Apostle said that widows living in luxury, though alive, were dead. Therefore, those are dead who do not have Christ living in them by faith and spirit, who do not have the virtue of Christ working in them—that is, who do not produce lively works. For the Lord is reported to have said also in the Gospel: “Let the dead bury their dead.” The Sardinians therefore had the name of living men, that is to say, they were called Christians, spiritual, regenerated, and holy worshippers of God; but they were dead, namely, hypocrites, in whom no spirit nor Christian life appeared. The flesh, the world, and corruption still lived in them. But such churches displease Christ. There are many such today. But how does Christ reject them? Truly, he condemns such, but not to destroy them—for so the world condemns—but that they should repent. For he does not will the death of a sinner, but rather that he should turn and live. And therefore, consequently, he prepares a medicine for the disease.

\index{Ezekiel (Book of)}First, he prescribes to the stars, or bishops, what they should do in this situation. Then, he tells the entire congregation their duty. From this, we learn how similar ailments in churches are to be healed. This belongs to the pastors, whom he commanded to watch diligently over the flock and to strengthen those who remained, not yet truly lost, but close to destruction, unless they are helped in time with sound and wholesome doctrine. He no doubt alluded to that pastoral care and responsibility which the Lord describes in chapter thirty-four of Ezekiel. The flock is strengthened by the word of God; by it, they are turned away from death and preserved in life, and so forth.

\index{Arethas of Caesarea}Now he also adds the reason why he commands that the flock be confirmed, lest they slip into death. For I have not found your works full or perfect before God. The Greek copies of Complutensian and Arethas read, “my God.” By works he understands all things that are done, words, deeds, and the entire conduct of people. The works, even those of the saints, are always imperfect if we consider human weakness. For as long as we live here, the flesh fights against the spirit; so that Job said he feared all his works and therefore fled to the clemency of the judge. Nevertheless, they are perfect and full in respect of Christ. For he is our fullness, and in him we are complete, as in John 1, Ephesians 1, and Colossians 1:12. And he makes us partakers of his fullness by faith. Those of Sardis were destitute of true faith, therefore every work of theirs must needs be imperfect before God, who allows nothing but that which comes from the Son and is most pure. Therefore the Lord commands that faith be taught diligently and instilled, that they may be made perfect in Christ. This is the best medicine for the deadly disease of Christ’s church.

\index{Repentance}Here follows the duty of the people, how they may be healed by apostolic repentance. The chief point is to remember the Lord’s words, in what we have heard and received them. We are not commanded to devise new forms of religion and repentance, but we are sent to the old tradition, not of men, but which we have in the Scriptures of the Evangelists and Apostles. These, I say, we ought to remember. For through custom of sinning, we forget God’s word. And truly, the beginning of Peter’s repentance was to have remembered the words of the Lord. Therefore, those who will not be reproved and instructed by God’s word shall never come to, or attain, true repentance.

Furthermore, it is necessary that we keep and retain the words of God, that is, the true doctrine of Christ, lest we forget it straightaway, or that we set it in vain contemplation and not in effectual work. The doctrine of Christ must be kept and performed in work. For in the last place it follows: and repent. True repentance consists in work: that in mind and body we should turn away from evil and turn unto God, and do good, being sorry for our wicked deeds past; this is the true apostolic repentance.

To which repentance, now, after the divine, prophetic, and apostolic manner, he draws people by means of threats. These threats are to be understood as applying both to ministers and to the people in the congregation. Again, the Lord uses parables, as we read in Matthew 24. With these, he exhorts to watchfulness and sobriety. Since that passage is explained at length, I need not use many words about it here. To the Lord be praise and thanks giving forevermore.
